\documentclass[11pt,letterpaper]{article}

% ============================================================
% COSMIC BREATH TIME CONVERTER IMPLEMENTATION MANUAL
% Compile with: pdfLaTeX
% Verified baseline date: July 19, 2026
% ============================================================

% ------------------------------------------------------------
% Core packages
% ------------------------------------------------------------
\usepackage[margin=0.78in,headheight=16pt]{geometry}
\usepackage[T1]{fontenc}
\usepackage[utf8]{inputenc}
\usepackage{lmodern}
\usepackage[scaled=0.94]{helvet}
\usepackage{microtype}
\usepackage{parskip}
\usepackage{setspace}
\usepackage{enumitem}
\usepackage{booktabs}
\usepackage{tabularx}
\usepackage{array}
\usepackage{longtable}
\usepackage{xcolor}
\usepackage{graphicx}
\usepackage{fancyhdr}
\usepackage{lastpage}
\usepackage{titlesec}
\usepackage{hyperref}
\usepackage{bookmark}
\usepackage{listings}
\usepackage[most]{tcolorbox}
\usepackage{amsmath}
\usepackage{amssymb}
\usepackage{pifont}

% ------------------------------------------------------------
% Document typography
% ------------------------------------------------------------
\renewcommand{\familydefault}{\sfdefault}
\setstretch{1.08}
\setlength{\parindent}{0pt}
\setlength{\parskip}{0.65em}

% ------------------------------------------------------------
% Cosmic Universalism color palette
% ------------------------------------------------------------
\definecolor{CUNavy}{HTML}{050812}
\definecolor{CUNavyTwo}{HTML}{080D1A}
\definecolor{CUSurface}{HTML}{0C1424}
\definecolor{CUGold}{HTML}{DDBF72}
\definecolor{CUGoldLight}{HTML}{FFF8DD}
\definecolor{CUCyan}{HTML}{35CDEB}
\definecolor{CUCyanLight}{HTML}{78E8FF}
\definecolor{CUBlue}{HTML}{668ED1}
\definecolor{CUBlueLight}{HTML}{9EB9EA}
\definecolor{CUAmber}{HTML}{D99A32}
\definecolor{CURed}{HTML}{B84C4C}
\definecolor{CUGreen}{HTML}{2C8A65}
\definecolor{CULight}{HTML}{F3F6FA}
\definecolor{CUBorder}{HTML}{CCD5E0}
\definecolor{CUText}{HTML}{18212E}
\definecolor{CUMuted}{HTML}{5C6878}
\definecolor{CUCodeBackground}{HTML}{F6F8FB}
\definecolor{CUCodeBorder}{HTML}{BFCADA}

% ------------------------------------------------------------
% Project metadata
% ------------------------------------------------------------
\newcommand{\ProjectTitle}{Cosmic Breath Time Converter}
\newcommand{\DocumentTitle}{Implementation Manual}
\newcommand{\RepositoryName}{CosmicUniversalismStatement}
\newcommand{\RepositoryURL}{https://github.com/willmaddock/CosmicUniversalismStatement}
\newcommand{\StableBranch}{website}
\newcommand{\WebsiteBranch}{feature/cu-time-converter}
\newcommand{\WebsiteDirectory}{website/}
\newcommand{\TargetRoute}{/CosmicUniversalismStatement/cu-time/}
\newcommand{\BaselineTag}{website-v1.1-cu-time-foundation}
\newcommand{\BaselineCommit}{4fd6897}
\newcommand{\DocumentVersion}{1.0}
\newcommand{\VerificationDate}{July 19, 2026}
\newcommand{\AuthorName}{William Maddock}

% ------------------------------------------------------------
% Hyperlinks and PDF metadata
% ------------------------------------------------------------
\hypersetup{
    colorlinks=true,
    linkcolor=CUBlue,
    urlcolor=CUBlue,
    citecolor=CUBlue,
    pdftitle={Cosmic Breath Time Converter Implementation Manual},
    pdfauthor={William Maddock},
    pdfsubject={Controlled migration of the Cosmic Breath Time Converter into Astro},
    pdfkeywords={Cosmic Universalism, CU-Time, Astro, TypeScript, Decimal.js,
                 Gregorian calendar, Julian Day Number, Vitest, functional parity}
}

% ------------------------------------------------------------
% Header and footer
% ------------------------------------------------------------
\pagestyle{fancy}
\fancyhf{}
\fancyhead[L]{\small\textbf{\ProjectTitle}}
\fancyhead[R]{\small\DocumentTitle}
\fancyfoot[L]{\small Version \DocumentVersion}
\fancyfoot[C]{\small \thepage\ of \pageref{LastPage}}
\fancyfoot[R]{\small Branch: \texttt{\WebsiteBranch}}
\renewcommand{\headrulewidth}{0.5pt}
\renewcommand{\footrulewidth}{0.5pt}

% ------------------------------------------------------------
% Section formatting
% ------------------------------------------------------------
\titleformat{\section}
  {\Large\bfseries\color{CUNavy}}
  {\thesection}
  {0.8em}
  {}
  [\vspace{0.2em}\color{CUCyan}\titlerule]

\titleformat{\subsection}
  {\large\bfseries\color{CUNavyTwo}}
  {\thesubsection}
  {0.7em}
  {}

\titleformat{\subsubsection}
  {\normalsize\bfseries\color{CUBlue}}
  {\thesubsubsection}
  {0.6em}
  {}

\titlespacing*{\section}{0pt}{2.2em}{1em}
\titlespacing*{\subsection}{0pt}{1.6em}{0.7em}
\titlespacing*{\subsubsection}{0pt}{1.2em}{0.4em}

% ------------------------------------------------------------
% Lists
% ------------------------------------------------------------
\setlist[itemize]{
    leftmargin=1.6em,
    itemsep=0.3em,
    topsep=0.3em
}

\setlist[enumerate]{
    leftmargin=2em,
    itemsep=0.45em,
    topsep=0.4em
}

\newcommand{\checkbox}{\(\square\)}
\newcommand{\checked}{\ding{51}}

% ------------------------------------------------------------
% Table helpers
% ------------------------------------------------------------
\newcolumntype{Y}{>{\raggedright\arraybackslash}X}
\newcolumntype{L}[1]{>{\raggedright\arraybackslash}p{#1}}

% ------------------------------------------------------------
% General information boxes
% ------------------------------------------------------------
\newtcolorbox{purposebox}{
    enhanced,
    breakable,
    colback=CULight,
    colframe=CUBlue,
    boxrule=0.8pt,
    arc=2mm,
    left=3mm,
    right=3mm,
    top=2.5mm,
    bottom=2.5mm,
    title=\textbf{Purpose},
    coltitle=white,
    colbacktitle=CUBlue
}

\newtcolorbox{checkpointbox}{
    enhanced,
    breakable,
    colback=CUCyan!6,
    colframe=CUCyan,
    boxrule=0.8pt,
    arc=2mm,
    left=3mm,
    right=3mm,
    top=2.5mm,
    bottom=2.5mm,
    title=\textbf{Checkpoint},
    coltitle=CUNavy,
    colbacktitle=CUCyanLight
}

\newtcolorbox{safetybox}{
    enhanced,
    breakable,
    colback=CUAmber!8,
    colframe=CUAmber,
    boxrule=0.8pt,
    arc=2mm,
    left=3mm,
    right=3mm,
    top=2.5mm,
    bottom=2.5mm,
    title=\textbf{Safety Rule},
    coltitle=white,
    colbacktitle=CUAmber
}

\newtcolorbox{warningbox}{
    enhanced,
    breakable,
    colback=CURed!6,
    colframe=CURed,
    boxrule=0.8pt,
    arc=2mm,
    left=3mm,
    right=3mm,
    top=2.5mm,
    bottom=2.5mm,
    title=\textbf{Do Not Continue Until Resolved},
    coltitle=white,
    colbacktitle=CURed
}

\newtcolorbox{successbox}{
    enhanced,
    breakable,
    colback=CUGreen!7,
    colframe=CUGreen,
    boxrule=0.8pt,
    arc=2mm,
    left=3mm,
    right=3mm,
    top=2.5mm,
    bottom=2.5mm,
    title=\textbf{Successful Result},
    coltitle=white,
    colbacktitle=CUGreen
}

\newtcolorbox{notebox}{
    enhanced,
    breakable,
    colback=CUCodeBackground,
    colframe=CUBorder,
    boxrule=0.7pt,
    arc=2mm,
    left=3mm,
    right=3mm,
    top=2.5mm,
    bottom=2.5mm
}

% ------------------------------------------------------------
% Listing styles
% ------------------------------------------------------------
\lstdefinestyle{terminalstyle}{
    basicstyle=\ttfamily\small,
    backgroundcolor=\color{CUNavyTwo},
    keywordstyle=\color{CUCyanLight},
    stringstyle=\color{CUGoldLight},
    commentstyle=\color{CUBlueLight},
    identifierstyle=\color{white},
    showstringspaces=false,
    columns=fullflexible,
    keepspaces=true,
    breaklines=true,
    breakatwhitespace=false,
    frame=none,
    tabsize=2,
    upquote=true
}

\lstdefinestyle{promptstyle}{
    basicstyle=\ttfamily\small,
    backgroundcolor=\color{CUCodeBackground},
    keywordstyle=\color{CUBlue},
    stringstyle=\color{CUGreen},
    commentstyle=\color{CUMuted},
    identifierstyle=\color{CUText},
    showstringspaces=false,
    columns=fullflexible,
    keepspaces=true,
    breaklines=true,
    breakatwhitespace=false,
    frame=none,
    tabsize=2,
    upquote=true
}

\lstdefinestyle{codestyle}{
    basicstyle=\ttfamily\small,
    backgroundcolor=\color{CUCodeBackground},
    keywordstyle=\color{CUBlue},
    stringstyle=\color{CUGreen},
    commentstyle=\color{CUMuted},
    identifierstyle=\color{CUText},
    showstringspaces=false,
    columns=fullflexible,
    keepspaces=true,
    breaklines=true,
    breakatwhitespace=false,
    frame=none,
    tabsize=2,
    upquote=true
}

% ------------------------------------------------------------
% Copy-and-paste boxes
% ------------------------------------------------------------
\newtcblisting{commandbox}[1]{
    enhanced,
    breakable,
    listing only,
    listing engine=listings,
    listing options={style=terminalstyle},
    colback=CUNavyTwo,
    colframe=CUCyan,
    coltitle=CUNavy,
    colbacktitle=CUCyanLight,
    boxrule=0.8pt,
    arc=2mm,
    left=2mm,
    right=2mm,
    top=1mm,
    bottom=1mm,
    title=\textbf{#1}
}

\newtcblisting{promptbox}[1]{
    enhanced,
    breakable,
    listing only,
    listing engine=listings,
    listing options={style=promptstyle},
    colback=CUCodeBackground,
    colframe=CUBlue,
    coltitle=white,
    colbacktitle=CUBlue,
    boxrule=0.8pt,
    arc=2mm,
    left=2mm,
    right=2mm,
    top=1mm,
    bottom=1mm,
    title=\textbf{Copy-and-Paste Prompt: #1}
}

\newtcblisting{codebox}[1]{
    enhanced,
    breakable,
    listing only,
    listing engine=listings,
    listing options={style=codestyle},
    colback=CUCodeBackground,
    colframe=CUCodeBorder,
    coltitle=white,
    colbacktitle=CUNavyTwo,
    boxrule=0.8pt,
    arc=2mm,
    left=2mm,
    right=2mm,
    top=1mm,
    bottom=1mm,
    title=\textbf{#1}
}

% ------------------------------------------------------------
% Convenience commands
% ------------------------------------------------------------
\newcommand{\filepath}[1]{\texttt{#1}}
\newcommand{\menu}[1]{\textbf{\textsf{#1}}}
\newcommand{\branchname}[1]{\texttt{#1}}
\newcommand{\statusline}[2]{\textbf{#1:} #2}

% ============================================================
% DOCUMENT
% ============================================================

\begin{document}

% ------------------------------------------------------------
% Cover page
% ------------------------------------------------------------
\begin{titlepage}
\pagecolor{CUNavy}
\color{white}
\thispagestyle{empty}

\vspace*{0.65in}

{\small\bfseries\color{CUCyanLight}
COSMIC UNIVERSALISM COMPUTATIONAL INTELLIGENCE INITIATIVE}

\vspace{0.42in}

{\Huge\bfseries
\ProjectTitle}

\vspace{0.18in}

{\LARGE\color{CUGoldLight}
\DocumentTitle}

\vspace{0.32in}

{\large
A specification-first, controlled migration manual for preserving the
approved converter behavior while integrating a testable TypeScript engine
and accessible interactive interface into the Astro CU-Time page.}

\vfill

\begin{tcolorbox}[
    colback=CUNavyTwo,
    colframe=CUBlue,
    boxrule=0.8pt,
    arc=2mm,
    width=\textwidth
]
\color{white}
\begin{tabularx}{\textwidth}{>{\bfseries}L{1.72in}Y}
Repository & \RepositoryName \\
Protected branch & \texttt{main} \\
Stable website branch & \texttt{\StableBranch} \\
Active feature branch & \texttt{\WebsiteBranch} \\
Foundation tag & \texttt{\BaselineTag} \\
Foundation commit & \texttt{\BaselineCommit} \\
Target route & \texttt{\TargetRoute} \\
Migration strategy & Approach B: tested TypeScript extraction \\
Reference implementation & \texttt{cosmic\_breath\_time\_converter\_v3\_0\_1.html} \\
Document version & \DocumentVersion \\
Verified & \VerificationDate \\
\end{tabularx}
\end{tcolorbox}

\vfill

{\large\bfseries \AuthorName}

\vspace{0.1in}

{\small \href{\RepositoryURL}{\RepositoryURL}}

\vspace{0.4in}

\end{titlepage}

\pagecolor{white}
\color{CUText}

% ------------------------------------------------------------
% Document control
% ------------------------------------------------------------
\section*{Document Control}
\addcontentsline{toc}{section}{Document Control}

\begin{tabularx}{\textwidth}{>{\bfseries}L{1.88in}Y}
\toprule
Document & \ProjectTitle: \DocumentTitle \\
Version & \DocumentVersion \\
Classification & Technical Specification and Controlled Migration Manual \\
Repository & \href{\RepositoryURL}{\RepositoryName} \\
Protected branch & \branchname{main} \\
Stable integration branch & \branchname{\StableBranch} \\
Active development branch & \branchname{\WebsiteBranch} \\
Foundation baseline & \branchname{\BaselineTag} at \texttt{\BaselineCommit} \\
Target page & \filepath{website/src/pages/cu-time.astro} \\
Target local route & \texttt{\TargetRoute} \\
Reference implementation & \filepath{cosmic\_breath\_time\_converter\_v3\_0\_1.html} \\
Technical baseline verified & \VerificationDate \\
\bottomrule
\end{tabularx}

\begin{purposebox}
This manual controls the solo-workflow migration of the existing standalone
Cosmic Breath Time Converter into the Cosmic Universalism Astro website.
It preserves approved numerical behavior while separating mathematical
specification, pure conversion logic, browser interaction, presentation,
testing, and release approval.
\end{purposebox}

\begin{safetybox}
This document does not approve any mathematical correction merely because a
potential defect or ambiguity is identified. The original HTML converter is
the behavioral reference. Every intentional change to constants, formulas,
rounding, calendar interpretation, labels, or outputs requires a documented
review decision and new test expectations.
\end{safetybox}

\clearpage
\tableofcontents
\clearpage

% ============================================================
\section*{START HERE - Initialize a New Project Chat}
\addcontentsline{toc}{section}{START HERE - Initialize a New Project Chat}
% ============================================================

Use this section whenever converter work is continued in a new chat inside
the existing \textbf{Cosmic Universalism Website} ChatGPT Project.

\subsection*{Required Project references}

Add or retain these files in the Project:

\begin{itemize}
    \item \filepath{Cosmic\_Universalism\_Website\_Implementation\_Manual\_v1.1.pdf}
    \item \filepath{Cosmic\_Breath\_Time\_Converter\_Implementation\_Manual\_v1.0.pdf}
    \item \filepath{cosmic-universalism-website-blueprint.md}
    \item \filepath{cosmic\_breath\_time\_converter\_v3\_0\_1.html}
\end{itemize}

\subsection*{New-chat bootstrap prompt}

\begin{promptbox}{Initialize the CU-Time Converter Chat}
We are continuing the controlled Cosmic Breath Time Converter
implementation inside the Cosmic Universalism Website Project.

Before proposing or making changes, read these Project references:

1. Cosmic_Breath_Time_Converter_Implementation_Manual_v1.0.pdf
2. Cosmic_Universalism_Website_Implementation_Manual_v1.1.pdf
3. cosmic_breath_time_converter_v3_0_1.html
4. cosmic-universalism-website-blueprint.md

Use the Project Instructions as mandatory governance rules.

Repository:
willmaddock/CosmicUniversalismStatement

Protected branch:
main

Stable website branch:
website

Active feature branch:
feature/cu-time-converter

Foundation baseline tag:
website-v1.1-cu-time-foundation

Foundation baseline commit:
4fd6897

Target page:
website/src/pages/cu-time.astro

Target route:
/CosmicUniversalismStatement/cu-time/

Reference implementation:
cosmic_breath_time_converter_v3_0_1.html

Selected migration strategy:
Approach B - preserve approved converter behavior while extracting
its mathematics into testable TypeScript modules and connecting an
accessible interface to the Astro CU-Time page.

The original HTML converter must remain unchanged as the reference
implementation.

First:

1. Summarize the applicable Project protection rules.
2. Summarize the verified foundation baseline.
3. Identify the current authorized manual phase.
4. Identify the target route and expected CU-Time files.
5. State which repository facts must be verified before editing.
6. State the next authorized task and its completion gate.

Do not edit files.
Do not install packages.
Do not switch branches.
Do not commit, tag, push, merge, or deploy.
Wait for repository verification before implementation.
\end{promptbox}

\begin{checkpointbox}
A new chat is initialized correctly only when it identifies
\branchname{main} as protected, \branchname{website} as the stable baseline,
\branchname{feature/cu-time-converter} as the active feature branch, and the
original HTML converter as an unchanged reference implementation.
\end{checkpointbox}

% ============================================================
\section{Part 1 - Project Governance and Verified Baseline}
% ============================================================

\subsection{Repository and branch authority}

\begin{tabularx}{\textwidth}{L{2.05in}Y}
\toprule
\textbf{Item} & \textbf{Verified state} \\
\midrule
Repository root & \nolinkurl{/Users/cosmos/Documents/GitHub/CosmicUniversalismStatement} \\
Origin & \href{\RepositoryURL}{\RepositoryURL.git} \\
Protected branch & \branchname{main}; never modify for this workflow \\
Stable website branch & \branchname{website} at \texttt{a4a4666} \\
Stable website tag & \branchname{website-v1.0-foundation} \\
Active feature branch & \branchname{feature/cu-time-converter} \\
Feature foundation commit & \nolinkurl{4fd68974e6a8ff99289087cb195a7b09f5759386} \\
Feature foundation tag & \branchname{website-v1.1-cu-time-foundation} \\
Working tree & Clean and tracking \texttt{origin/feature/cu-time-converter} \\
Deployment trigger & \branchname{website} only \\
\bottomrule
\end{tabularx}

\subsection{Verified technology state}

\begin{tabularx}{\textwidth}{L{2.05in}Y}
\toprule
\textbf{Technology} & \textbf{Verified state} \\
\midrule
Node.js & \texttt{v24.18.0} \\
npm & \texttt{11.16.0} \\
Astro & \texttt{\string^7.1.1} \\
TypeScript configuration & Astro strict configuration \\
Astro site & \texttt{https://willmaddock.github.io} \\
Astro base & \texttt{/CosmicUniversalismStatement} \\
Arbitrary-precision dependency & Not installed \\
Automated test framework & Not installed \\
Existing client scripts & None detected in \filepath{website/src/} \\
\bottomrule
\end{tabularx}

\subsection{Verified foundation files}

\begin{codebox}{Existing CU-Time files}
website/src/pages/cu-time.astro
website/src/components/CUTimeConverter.astro
website/src/components/CUTimeInvitation.astro
website/src/lib/cuTime/constants.ts
\end{codebox}

The current page is explanatory and imports the foundation component. The
component contains accessible presentation markup, classification, and an
active-development status. It does not contain controls or calculations.

The constants workspace is not empty. It contains an initial inventory of
NASA-aligned and CU framework values as ordinary JavaScript number literals.
Those values are not yet approved for calculation and may not preserve all
source digits when evaluated as binary floating-point numbers.

\begin{warningbox}
Do not describe \filepath{constants.ts} as an empty placeholder. The accurate
status is: \textbf{inventory extracted, precision representation and
mathematical approval pending}.
\end{warningbox}

\subsection{Verified production build}

The foundation build was executed on \VerificationDate.

\begin{codebox}{Verified build result}
Command: npm run build
Exit code: 0
Output mode: static
Pages built: 7
CU-Time route generated: /cu-time/index.html
Active branch: feature/cu-time-converter
HEAD: 4fd6897
Working tree after build: clean
\end{codebox}

\begin{successbox}
The CU-Time presentation foundation is buildable and the target route exists.
This authorizes specification work. It does not yet authorize mathematical
integration.
\end{successbox}

\subsection{Project Instructions appendix block}

Append this block to the existing ChatGPT Project Instructions. Do not replace
the original website instructions.

\begin{promptbox}{Permanent CU-Time Project Instructions}
CURRENT CU-TIME DEVELOPMENT WORKFLOW

ACTIVE FEATURE:
Cosmic Breath Time Converter integration

STABLE WEBSITE BRANCH:
website

STABLE WEBSITE MILESTONE:
website-v1.0-foundation

STABLE WEBSITE COMMIT:
a4a4666

ACTIVE DEVELOPMENT BRANCH:
feature/cu-time-converter

CURRENT FEATURE FOUNDATION:
website-v1.1-cu-time-foundation

FOUNDATION COMMIT:
4fd6897

TARGET PAGE:
website/src/pages/cu-time.astro

TARGET LOCAL ROUTE:
/CosmicUniversalismStatement/cu-time/

CURRENT FOUNDATION COMPONENT:
website/src/components/CUTimeConverter.astro

CURRENT CONSTANTS INVENTORY:
website/src/lib/cuTime/constants.ts

REFERENCE CONVERTER:
cosmic_breath_time_converter_v3_0_1.html

IMPLEMENTATION APPROACH:
Approach B - extract approved converter behavior into pure,
testable TypeScript modules and integrate it with the Astro CU-Time
page using lightweight browser JavaScript.

FUNCTIONAL PARITY RULE:
The integrated converter must preserve approved numerical behavior,
precision, inputs, outputs, calendar handling, and conversion
capabilities of the reference converter. The website interface does
not need to preserve the reference converter's React, CDN, Tailwind,
or visual architecture.

REFERENCE PROTECTION:
Do not modify the original HTML reference converter during the
initial migration.

MATHEMATICAL PROTECTION:
1. Do not silently change constants.
2. Do not silently change formulas.
3. Do not silently change rounding behavior.
4. Do not silently change calendar assumptions.
5. Record each proposed correction separately.
6. Establish golden reference vectors before implementing the new
   conversion engine.
7. Demonstrate parity before introducing enhancements.

BRANCH PROTECTION:
1. Never modify main.
2. Keep website as the stable website baseline.
3. Perform CU-Time work only on feature/cu-time-converter.
4. Confirm branch and working-tree status before every edit.
5. Do not merge into website without explicit approval.
6. Do not commit, tag, push, merge, or deploy unless explicitly
   requested.

CURRENT DEVELOPMENT BOUNDARY:
The presentation foundation is complete. Constants review,
mathematical specification, engine implementation, automated tests,
and interactive integration remain pending.
\end{promptbox}

\subsection{Required pre-edit terminal check}

\begin{commandbox}{Verify the live repository before every editing session}
cd "$(git rev-parse --show-toplevel)" || exit 1

git branch --show-current
git rev-parse --short HEAD
git status --short --branch
\end{commandbox}

Expected branch:

\begin{codebox}{Required branch}
feature/cu-time-converter
\end{codebox}

\subsection{No-edit Codex inspection prompt}

\begin{promptbox}{Verify the Live Repository}
Inspect the current repository without changing any files.

Expected active branch:
feature/cu-time-converter

Foundation baseline tag:
website-v1.1-cu-time-foundation

Foundation baseline commit:
4fd6897

Verify and report:

1. Repository root and origin
2. Active branch and HEAD
3. Working-tree status
4. Relevant website and CU-Time tags
5. CU-Time-related file structure
6. Astro configuration and base-path handling
7. package.json scripts and dependencies
8. Existing test infrastructure
9. Existing arbitrary-precision dependency
10. Current CU-Time page and component responsibilities
11. Current constants inventory
12. Existing client-script conventions
13. Deployment workflow trigger branch
14. Production build status
15. Differences from this manual

Do not install packages.
Do not create or edit files.
Do not switch branches.
Do not commit, tag, push, merge, or deploy.
Return a structured verification report only.
\end{promptbox}

% ============================================================
\section{Part 2 - Scope and Functional Parity}
% ============================================================

\subsection{Migration decision}

The approved migration is \textbf{Approach B}: retain the standalone HTML
converter unchanged as a reference and extract approved behavior into pure
TypeScript modules. The Astro website provides the integrated presentation and
browser interaction.

\begin{tabularx}{\textwidth}{L{2.25in}Y}
\toprule
\textbf{Parity category} & \textbf{Requirement} \\
\midrule
Mathematical behavior & Required to match approved reference outputs \\
Decimal precision & Required; exact literals must enter Decimal from strings \\
Forward conversion & Required \\
Reverse conversion & Required \\
CE and BCE support & Required after calendar specification approval \\
NASA-aligned output & Required \\
Years-since-Big-Bang output & Required \\
Essential validation & Required \\
Copy and clear controls & Required for first interactive release \\
React architecture & Not required \\
Tailwind classes & Not required \\
CDN dependencies & Not permitted for the calculation engine \\
Original visual design & Not required \\
\bottomrule
\end{tabularx}

\subsection{First interactive release scope}

\textbf{Included:}

\begin{itemize}
    \item Gregorian date and time to CU-Time.
    \item CU-Time to Gregorian date and time.
    \item CE/BCE selection and documented astronomical-year mapping.
    \item NASA-aligned CU-Time.
    \item Years since Big Bang as defined by the approved converter model.
    \item Full numeric and scientific-notation results.
    \item Accessible validation and result announcements.
    \item Copy and clear controls.
    \item Methodology and classification disclosure.
\end{itemize}

\textbf{Deferred until core parity is approved:}

\begin{itemize}
    \item CSV upload and export.
    \item XLSX support.
    \item Embedded Markdown research and About tabs.
    \item CU-Time arithmetic calculator.
    \item Theme switching.
    \item Spreadsheet editing.
    \item Additional cosmological phase calculations.
\end{itemize}

\subsection{Reference behavior inventory prompt}

\begin{promptbox}{Audit the Reference Converter Without Editing}
Read the complete reference file:
cosmic_breath_time_converter_v3_0_1.html

Do not modify it.

Produce a behavioral inventory with these sections:

1. External runtime dependencies and versions
2. Decimal precision and rounding configuration
3. Constants used directly by conversion functions
4. Constants mentioned only in research text
5. Input fields and validation behavior
6. Gregorian-to-CU conversion path
7. CU-to-Gregorian conversion path
8. NASA-aligned calculations
9. Formatting functions
10. Metadata calculations
11. Result fields and labels
12. Calculator behavior
13. CSV/XLSX behavior
14. Accessibility behavior
15. Potential ambiguities or defects
16. Features required for first-release parity
17. Features safe to defer

For every potential defect, label it REVIEW REQUIRED. Do not
silently propose a corrected expected value.

Do not edit repository files.
Do not install packages.
Do not commit, tag, push, merge, or deploy.
\end{promptbox}

\subsection{Parity decision record}

Every behavior must receive one decision:

\begin{codebox}{Permitted parity decisions}
PRESERVE EXACTLY
PRESERVE WITH DOCUMENTED REPRESENTATION CHANGE
CORRECT AFTER APPROVAL
DEFER FROM FIRST RELEASE
REMOVE AFTER APPROVAL
OPEN QUESTION
\end{codebox}

\begin{checkpointbox}
No engine implementation begins until the reference behavior inventory and
parity decision record are reviewed and accepted.
\end{checkpointbox}

% ============================================================
\section{Part 3 - Mathematical Constants}
% ============================================================

\subsection{Direct reference-engine constants}

The standalone converter constructs these values with Decimal strings and sets
precision to 60 digits with half-up rounding.

\begin{tabularx}{\textwidth}{L{1.7in}Y L{1.35in}}
\toprule
\textbf{Name} & \textbf{Exact source literal} & \textbf{Initial status} \\
\midrule
\texttt{ANCHOR\_YEAR} & \nolinkurl{2000} & Review \\
\texttt{ANCHOR\_JDN} & \nolinkurl{2451544.5} & Review \\
\texttt{DAYS\_PER\_YEAR} & \nolinkurl{365.2421897} & Review \\
\texttt{BASE\_CU\_NASA} & \nolinkurl{3094213000000} & Review \\
\texttt{OFFSET} & \nolinkurl{78955076.49024643522707769749119} & Precision critical \\
\texttt{NASA\_UNIVERSE\_AGE} & \nolinkurl{13786999981.453} & Review \\
\texttt{BIG\_BANG\_CU\_NASA} & Derived by subtraction & Review \\
\bottomrule
\end{tabularx}

\subsection{Current website constants inventory}

The existing \nolinkurl{website/src/lib/cuTime/constants.ts} also lists these
CU framework values:

\begin{codebox}{Additional constants currently inventoried}
baseCu
subZtomCu
secondsPerDay
secondsPerYear
cosmicLifespan
convergenceYear
ctomStart
ztomProximity
\end{codebox}

They require separate provenance review because not all are used by the
reference engine's core conversion functions.

\begin{warningbox}
Long decimal constants must not remain ordinary JavaScript number literals in
the approved engine. Binary floating-point evaluation can discard source
digits before Decimal receives the value. Store exact values as strings and
construct Decimal instances from those strings.
\end{warningbox}

\subsection{Proposed precision-safe representation}

\begin{codebox}{Illustrative constants representation - not yet authorized code}
export const CU_TIME_CONSTANT_LITERALS = {
  anchorYear: '2000',
  anchorJdn: '2451544.5',
  daysPerYear: '365.2421897',
  baseCuNasa: '3094213000000',
  offset: '78955076.49024643522707769749119',
  nasaUniverseAge: '13786999981.453',
} as const;
\end{codebox}

\subsection{Constant-review worksheet}

For each constant, record:

\begin{longtable}{L{1.5in}L{4.85in}}
\toprule
\textbf{Field} & \textbf{Required entry} \\
\midrule
Name & Exact code identifier \\
Source literal & Character-for-character source value \\
Units & Years, days, JDN, CU-Time, seconds, or dimensionless \\
Meaning & What the constant represents \\
Provenance & Source file and research reference \\
Used by & Exact functions or equations \\
Precision & Required significant and decimal digits \\
Uncertainty & Empirical, model-defined, or not specified \\
Approval & Approved, provisional, rejected, or open \\
Change control & Whether parity expectation changes \\
\bottomrule
\end{longtable}

\subsection{Constant audit prompt}

\begin{promptbox}{Create the Constant Registry}
Create a proposed CU-Time constant registry from:

- cosmic_breath_time_converter_v3_0_1.html
- website/src/lib/cuTime/constants.ts
- approved CU-Time research references

Do not edit files.

Separate constants into:

A. Direct conversion-engine constants
B. Derived constants
C. Research-only constants
D. UI-only values
E. Unresolved or conflicting constants

For every constant, document:

- exact identifier
- exact literal as text
- units
- meaning
- source location
- functions that use it
- precision requirement
- uncertainty or model status
- whether it is approved for the first engine

Detect values that would lose precision as JavaScript numbers.
Do not change values.
Mark discrepancies as REVIEW REQUIRED.
\end{promptbox}

\begin{checkpointbox}
Part 3 is complete only after the direct engine constants are approved as
exact strings, derived values have explicit formulas, and research-only values
are excluded from the first conversion engine unless separately authorized.
\end{checkpointbox}

% ============================================================
\section{Part 4 - Calendar and Time Specification}
% ============================================================

\subsection{Reference calendar behavior}

The reference converter uses a proleptic Gregorian-style Julian Day Number
algorithm. It maps BCE input to astronomical years with:

\[
Y_{\mathrm{astronomical}} = 1 - Y_{\mathrm{BCE}}.
\]

Therefore, 1 BCE maps to astronomical year 0, and 2 BCE maps to astronomical
year -1.

The anchor JDN is \texttt{2451544.5}. The implementation treats the input time
as UTC-like civil time and adds a day fraction based on 86,400 seconds.

\subsection{Specification decisions required}

\begin{itemize}
    \item Confirm the proleptic Gregorian calendar across all supported years.
    \item Confirm whether year 0 is accepted as direct CE input or only used internally.
    \item Confirm BCE display and astronomical-year conversion.
    \item Confirm the midnight/noon interpretation of the JDN anchor.
    \item Confirm UTC as the only supported time basis.
    \item Define valid month/day combinations and leap-year rules.
    \item Define maximum supported year and safe conversion limits.
    \item Define second precision and fractional-second policy.
    \item Define rollover behavior when rounded seconds reach 60.
\end{itemize}

\subsection{Gregorian leap-year rule}

For the proleptic Gregorian calendar, a year is a leap year when it is
divisible by 4, except years divisible by 100 are not leap years unless they
are divisible by 400. BCE application must use the approved astronomical-year
mapping consistently.

\subsection{Reference JDN formula}

For year $Y$, month $M$, day $D$, and time values $h$, $m$, and $s$:

\begin{align*}
a &= \left\lfloor\frac{14-M}{12}\right\rfloor, \\
y &= Y + 4800 - a, \\
m' &= M + 12a - 3, \\
J_{\mathrm{int}} &= D + \left\lfloor\frac{153m'+2}{5}\right\rfloor
+ 365y + \left\lfloor\frac{y}{4}\right\rfloor
- \left\lfloor\frac{y}{100}\right\rfloor
+ \left\lfloor\frac{y}{400}\right\rfloor - 32045, \\
f &= \frac{3600h+60m+s}{86400}, \\
J &= J_{\mathrm{int}} - 0.5 + f.
\end{align*}

This records the reference algorithm. It is not a scientific endorsement of
all supported date extrapolations.

\subsection{Known validation gap}

The reference input validation checks ranges such as month 1--12 and day
1--31, but does not fully reject invalid combinations such as April 31 or a
non-leap-year February 29 before conversion.

\begin{warningbox}
Improved date validation is a desirable correction, but it changes observable
behavior. Record it as a separately approved correction with explicit tests.
\end{warningbox}

\subsection{Calendar specification prompt}

\begin{promptbox}{Draft the Calendar Specification}
Draft the formal calendar and time specification for the integrated
converter.

Do not edit code.

Address:

1. Proleptic Gregorian calendar definition
2. Julian Day Number convention
3. Anchor JDN and its civil date/time meaning
4. UTC handling
5. CE/BCE input and output
6. Astronomical year zero mapping
7. Leap-year rules
8. Month-length validation
9. Maximum supported year
10. Seconds and fractional seconds
11. Rounding and day rollover
12. Invalid-input behavior
13. Round-trip tolerance
14. Known differences between the reference behavior and the
    proposed validated behavior

Label each proposed behavioral change REVIEW REQUIRED.
Do not silently approve corrections.
\end{promptbox}

% ============================================================
\section{Part 5 - Gregorian-to-CU Conversion}
% ============================================================

\subsection{Reference conversion path}

The approved path to be specified and tested is:

\begin{codebox}{Forward conversion sequence}
Gregorian date, time, and era
  -> astronomical year
  -> Julian Day Number
  -> delta days from ANCHOR_JDN
  -> delta years using DAYS_PER_YEAR
  -> NASA-aligned CU-Time
  -> subtract OFFSET
  -> converter-aligned CU-Time
  -> derived and formatted outputs
\end{codebox}

\subsection{Reference equations}

Let $J$ be the calculated JDN and $J_0$ the anchor JDN.

\begin{align*}
\Delta J &= J - J_0, \\
\Delta Y &= \frac{\Delta J}{D_Y}, \\
CU_{\mathrm{NASA}} &= CU_{\mathrm{NASA},0} + \Delta Y, \\
CU_{\mathrm{converter}} &= CU_{\mathrm{NASA}} - O, \\
Y_{\mathrm{BigBang}} &= Y_{\mathrm{NASA},0}
  + \left(CU_{\mathrm{NASA}} - CU_{\mathrm{NASA},0}\right),
\end{align*}

where $D_Y$ is \texttt{DAYS\_PER\_YEAR} and $O$ is \texttt{OFFSET}.

\subsection{Required forward output fields}

\begin{itemize}
    \item Original input date, era, and time.
    \item Normalized Gregorian date.
    \item Converter-aligned CU-Time formatted by magnitude.
    \item Full converter-aligned numeric value.
    \item Scientific notation.
    \item NASA-aligned CU-Time formatted by magnitude.
    \item Full NASA-aligned numeric value.
    \item NASA-aligned scientific notation.
    \item Years since Big Bang according to the converter model.
    \item Approved metadata fields.
\end{itemize}

\subsection{Potential formatting ambiguity}

The reference \texttt{formatCUTime} function appears to combine a remainder
that may already contain the fractional component with \texttt{cuTime.mod(1)}.
This may double-count part of the fraction. The behavior must be captured in a
golden vector before any correction is considered.

\subsection{Forward specification prompt}

\begin{promptbox}{Specify Gregorian-to-CU Conversion}
Using the approved constants registry and calendar specification,
draft the normative Gregorian-to-CU conversion specification.

Do not implement code.

Include:

- typed inputs
- normalization rules
- validation order
- exact Decimal operations
- forward equations
- rounding points
- result object fields
- display formatting
- error conditions
- test invariants
- reference-converter comparison procedure

Identify any reference behavior that appears ambiguous or defective.
Label it REVIEW REQUIRED and preserve the original expected output
until an explicit correction is approved.
\end{promptbox}

\begin{checkpointbox}
The forward algorithm is authorized for implementation only after its inputs,
equations, output schema, rounding behavior, and golden expectations are
approved.
\end{checkpointbox}

% ============================================================
\section{Part 6 - CU-to-Gregorian Conversion}
% ============================================================

\subsection{Reference reverse path}

\begin{codebox}{Reverse conversion sequence}
Converter-aligned CU-Time
  -> add OFFSET
  -> NASA-aligned CU-Time
  -> delta years from BASE_CU_NASA
  -> Julian Day Number
  -> Gregorian calendar conversion
  -> CE/BCE display
  -> derived and formatted outputs
\end{codebox}

\subsection{Reference reverse equations}

\begin{align*}
CU_{\mathrm{NASA}} &= CU_{\mathrm{converter}} + O, \\
\Delta Y &= CU_{\mathrm{NASA}} - CU_{\mathrm{NASA},0}, \\
J &= J_0 + \Delta Y D_Y.
\end{align*}

The JDN is then converted to Gregorian components using the reference
\texttt{jdnToDate} algorithm.

\subsection{Reverse-conversion review points}

\begin{itemize}
    \item Decimal input syntax and whitespace handling.
    \item Negative CU-Time support.
    \item Maximum supported magnitude.
    \item Gregorian year conversion from Decimal to JavaScript number.
    \item Rounding seconds to the nearest integer.
    \item Carrying seconds to minutes and minutes to hours.
    \item Carrying hour 24 into the next calendar day.
    \item CE/BCE display convention.
    \item Round-trip comparison tolerance.
\end{itemize}

The reference function reduces hours from 24 to 0 but does not visibly
increment the calendar date in that branch. This is a potential rollover defect
and must be evaluated with a targeted vector.

\subsection{Reverse specification prompt}

\begin{promptbox}{Specify CU-to-Gregorian Conversion}
Using the approved constants registry and calendar specification,
draft the normative CU-to-Gregorian conversion specification.

Do not implement code.

Include:

- accepted Decimal input grammar
- exact offset and JDN equations
- Gregorian reconstruction algorithm
- CE/BCE display rules
- rounding and carry rules
- next-day rollover behavior
- output object fields
- error handling
- maximum supported magnitude
- round-trip invariants
- reference parity comparison

Create targeted REVIEW REQUIRED cases for second, minute, hour,
day, month, year, leap-day, and BCE rollovers.
\end{promptbox}

% ============================================================
\section{Part 7 - Golden Reference Test Vectors}
% ============================================================

\subsection{Purpose}

A golden vector records an input and the complete output produced by the
unchanged reference converter. Golden vectors protect behavior during the
TypeScript migration. They do not automatically prove that the underlying
mathematics is correct.

\subsection{Required vector set}

\begin{longtable}{L{0.55in}L{2.65in}L{2.9in}}
\toprule
\textbf{ID} & \textbf{Input} & \textbf{Reason} \\
\midrule
F-01 & 01/01/2000 CE 00:00:00 & Anchor date \\
F-02 & 06/05/2025 CE 00:12:00 & Metadata reference instant \\
F-03 & 06/05/2025 CE 00:00:00 & Nearby reference boundary \\
F-04 & 07/19/2026 CE 00:00:00 & Manual verification date \\
F-05 & 02/29/2000 CE 12:34:56 & Valid leap day \\
F-06 & 02/29/1900 CE 00:00:00 & Invalid Gregorian leap day \\
F-07 & 04/31/2025 CE 00:00:00 & Invalid month/day combination \\
F-08 & 01/01/0001 CE 00:00:00 & CE boundary \\
F-09 & 01/01/0001 BCE 00:00:00 & Astronomical year zero mapping \\
F-10 & Maximum supported year & Magnitude boundary \\
R-01 & Anchor CU-Time & Reverse anchor \\
R-02 & Approved present CU-Time & Reference example \\
R-03 & Negative CU-Time & Input domain \\
R-04 & Rollover-sensitive CU-Time & Seconds/day carry \\
RT-01 & Forward then reverse & Round-trip invariant \\
RT-02 & Reverse then forward & Round-trip invariant \\
\bottomrule
\end{longtable}

\subsection{Derived anchor candidate}

At the anchor JDN, the reference equations give:

\[
CU_{\mathrm{converter},0}
= 3094213000000
- 78955076.49024643522707769749119.
\]

This produces the source literal represented in the current website inventory
as approximately \nolinkurl{3094134044923.509753564772922302508810}. It must still
be captured from the running reference converter and compared character by
character before becoming a golden value.

\subsection{Golden vector record format}

\begin{codebox}{Required vector record}
ID:
Reference converter version:
Reference file hash:
Input mode:
Input date/time/era or CU-Time:
Observed output labels:
Observed full numeric values:
Observed formatted values:
Observed scientific notation:
Observed metadata:
Browser and version:
Capture date:
Reviewer:
Parity decision:
Notes and open questions:
\end{codebox}

\subsection{Reference hash command}

\begin{commandbox}{Record the unchanged reference file hash}
shasum -a 256 cosmic_breath_time_converter_v3_0_1.html
\end{commandbox}

Use the actual repository path when the file is stored elsewhere.

\subsection{Golden vector prompt}

\begin{promptbox}{Prepare the Golden Vector Plan}
Create the golden reference vector worksheet for the unchanged
HTML converter.

Do not edit the reference converter or website code.

For each required vector:

- define exact input
- define which output fields must be captured
- identify the behavior being tested
- identify expected validation behavior
- identify whether the case is parity-only or correction-review
- define a reproducible capture procedure

Do not invent outputs that have not been observed from the reference
converter. Leave observed-output fields blank until captured.
\end{promptbox}

\subsection{Golden vector approval gate}

\begin{itemize}
    \item[\checkbox] Reference file hash recorded.
    \item[\checkbox] All required forward vectors captured.
    \item[\checkbox] All required reverse vectors captured.
    \item[\checkbox] Invalid-date behavior captured.
    \item[\checkbox] BCE behavior captured.
    \item[\checkbox] Rollover behavior captured.
    \item[\checkbox] Full precision copied without spreadsheet conversion.
    \item[\checkbox] Each ambiguity has a review decision.
    \item[\checkbox] Golden vector dataset approved before engine coding.
\end{itemize}

% ============================================================
\section{Part 8 - TypeScript Conversion Engine Architecture}
% ============================================================

\subsection{Approved architecture direction}

The engine should contain pure modules with no direct DOM access. The Astro
component and browser controller should consume the engine through typed
interfaces.

\begin{codebox}{Target CU-Time architecture}
website/src/lib/cuTime/
|-- constants.ts
|-- types.ts
|-- decimal.ts
|-- validation.ts
|-- calendar.ts
|-- converter.ts
|-- format.ts
|-- metadata.ts
`-- index.ts

website/src/components/
`-- CUTimeConverter.astro

website/src/scripts/
`-- cu-time-converter.ts

website/src/pages/
`-- cu-time.astro

website/src/lib/cuTime/__tests__/
|-- constants.test.ts
|-- calendar.test.ts
|-- converter.test.ts
|-- format.test.ts
`-- golden-vectors.test.ts
\end{codebox}

\subsection{Module responsibilities}

\begin{tabularx}{\textwidth}{L{1.65in}Y}
\toprule
\textbf{Module} & \textbf{Responsibility} \\
\midrule
\texttt{constants.ts} & Exact approved literal registry; no calculations \\
\texttt{types.ts} & Public inputs, outputs, errors, era, and vector types \\
\texttt{decimal.ts} & Decimal configuration and safe constructors \\
\texttt{validation.ts} & Date, time, era, numeric, and domain validation \\
\texttt{calendar.ts} & Gregorian/JDN conversion only \\
\texttt{converter.ts} & Forward, reverse, NASA, and Big-Bang calculations \\
\texttt{format.ts} & Numeric and magnitude presentation \\
\texttt{metadata.ts} & Approved secondary result calculations \\
\texttt{index.ts} & Controlled public exports \\
Browser controller & Form events, state, rendering, copy, and clear \\
Astro component & Semantic form, results, labels, status, and styles \\
\bottomrule
\end{tabularx}

\subsection{Dependency plan}

No arbitrary-precision dependency or test framework is currently installed.
After specification approval, the intended dependency additions are:

\begin{commandbox}{Install approved dependencies - run only when authorized}
cd "$(git rev-parse --show-toplevel)/website" || exit 1

npm install decimal.js
npm install --save-dev vitest
\end{commandbox}

The exact package versions written to \filepath{package-lock.json} become part
of the milestone record.

\subsection{Required package scripts}

\begin{codebox}{Proposed package.json scripts}
"test": "vitest run",
"test:watch": "vitest",
"check": "astro check"
\end{codebox}

If \texttt{astro check} is added, the compatible checker package must also be
installed according to the then-current Astro requirements.

\subsection{Type and error design}

The engine should return structured data and typed errors rather than formatted
HTML strings. Suggested categories include:

\begin{codebox}{Suggested error codes}
INVALID_DATE_FORMAT
INVALID_TIME_FORMAT
INVALID_MONTH
INVALID_DAY
INVALID_LEAP_DAY
INVALID_ERA
INVALID_CU_TIME
OUT_OF_SUPPORTED_RANGE
DIVISION_BY_ZERO
PRECISION_LIMIT
INTERNAL_CONVERSION_ERROR
\end{codebox}

\subsection{Architecture planning prompt}

\begin{promptbox}{Finalize the TypeScript Architecture}
Using the approved specification and verified Astro repository,
finalize the TypeScript engine architecture.

Do not edit files.

For every proposed file, define:

- responsibility
- imports
- exports
- public types
- error behavior
- Decimal boundary
- test ownership
- prohibited responsibilities

Requirements:

- pure conversion functions
- exact string literals for precision-critical constants
- no DOM access in lib/cuTime
- no React
- no Tailwind dependency
- no calculation CDN
- lightweight browser controller
- compatibility with Astro strict TypeScript
- compatibility with the configured base path
- behavior controlled by golden vectors

Return the file plan and dependency plan only.
\end{promptbox}

\begin{checkpointbox}
The architecture is complete only when each responsibility has one owner,
calculation code is independent of the UI, and no precision-critical value
enters Decimal through a JavaScript number.
\end{checkpointbox}

% ============================================================
\section{Part 9 - Automated Testing}
% ============================================================

\subsection{Testing layers}

\begin{enumerate}
    \item \textbf{Constant tests}: exact literals and derived identities.
    \item \textbf{Calendar tests}: JDN, leap years, BCE, and rollovers.
    \item \textbf{Conversion tests}: forward and reverse equations.
    \item \textbf{Golden tests}: complete parity with captured reference outputs.
    \item \textbf{Round-trip tests}: input recovery within approved tolerance.
    \item \textbf{Validation tests}: malformed and impossible inputs.
    \item \textbf{Formatting tests}: magnitude labels and full numeric output.
    \item \textbf{Build test}: Astro static generation.
    \item \textbf{Manual accessibility test}: keyboard and screen-reader behavior.
\end{enumerate}

\subsection{Test commands}

\begin{commandbox}{Run the automated and production checks}
cd "$(git rev-parse --show-toplevel)/website" || exit 1

npm test
npm run build
\end{commandbox}

\subsection{Numeric comparison policy}

Use exact string equality when the reference output is specified as exact and
both implementations use the same approved Decimal operations. Use an explicit
Decimal tolerance only when the specification documents a rounding boundary or
conversion loss. Never use an undocumented floating-point epsilon.

\subsection{Minimum test requirements}

\begin{itemize}
    \item[\checkbox] Anchor JDN conversion.
    \item[\checkbox] Anchor CU-Time conversion.
    \item[\checkbox] All golden vectors.
    \item[\checkbox] CE/BCE boundary.
    \item[\checkbox] Leap-year and invalid leap-day behavior.
    \item[\checkbox] Month-length validation.
    \item[\checkbox] Midnight and near-midnight rollover.
    \item[\checkbox] Negative and large CU-Time inputs.
    \item[\checkbox] Invalid numeric syntax.
    \item[\checkbox] Full precision literal preservation.
    \item[\checkbox] Format behavior, including the fractional remainder review.
    \item[\checkbox] Metadata reference instant and units.
    \item[\checkbox] Round-trip invariants.
\end{itemize}

\subsection{Test implementation prompt}

\begin{promptbox}{Implement the Approved Test Harness}
Implement only the approved automated test foundation and test
vectors.

Preconditions:

- active branch is feature/cu-time-converter
- working tree is understood
- constants registry is approved
- calendar specification is approved
- golden vectors are approved

Requirements:

- install only approved dependencies
- configure Vitest minimally
- add test scripts
- store golden vectors in a reviewable text-based format
- test exact strings where required
- use documented Decimal tolerances only where approved
- do not implement the interactive UI
- do not modify the reference HTML converter
- run npm test
- run npm run build
- report every changed file and dependency

Do not commit, tag, push, merge, or deploy.
\end{promptbox}

\subsection{Testing completion gate}

\begin{successbox}
The engine may be connected to the page only after all approved unit and golden
vector tests pass and the production build succeeds.
\end{successbox}

% ============================================================
\section{Part 10 - Interactive Astro Page Integration}
% ============================================================

\subsection{Integration objective}

Replace the current ``Interactive converter under development'' content with
an accessible working interface while preserving the page's documentation,
classification, research-status language, and visual system.

\subsection{First-release interface}

\begin{itemize}
    \item Mode selector or clearly separated forward and reverse forms.
    \item Date input with explicit \texttt{MM/DD/YYYY} help.
    \item Time input with explicit 24-hour \texttt{HH:MM:SS} help.
    \item CE/BCE selector.
    \item CU-Time decimal input.
    \item Convert controls.
    \item Clear controls.
    \item Copy result control.
    \item Validation summary and field-level errors.
    \item Result region announced with \texttt{aria-live}.
    \item Full precision result preserved as text.
    \item Scientific and formatted display values.
    \item Calculation-method disclosure.
\end{itemize}

\subsection{Astro client strategy}

The project currently contains no detected client scripts and no UI framework.
The preferred first implementation is:

\begin{codebox}{Client strategy}
Astro semantic markup
  + processed TypeScript <script>
  + imported pure CU-Time modules
  + DOM event listeners scoped to the component
  + no React island
  + no browser-side Babel
  + no Tailwind CDN
\end{codebox}

The browser controller should progressively enhance valid HTML. All calculation
logic remains in \nolinkurl{website/src/lib/cuTime/}.

\subsection{Route and base-path requirements}

The source page remains:

\begin{codebox}{Source and route}
Source: website/src/pages/cu-time.astro
Development/public route:
/CosmicUniversalismStatement/cu-time/
Generated static file:
website/dist/cu-time/index.html
\end{codebox}

Use the existing centralized base-path handling for internal links and assets.
The calculator itself should not hard-code the repository base into mathematical
logic.

\subsection{Accessibility requirements}

\begin{itemize}
    \item Every field has a visible label.
    \item Format help is associated with the field.
    \item Errors identify the field and corrective action.
    \item Focus moves only when necessary and predictably.
    \item Results are announced without stealing keyboard focus.
    \item Copy buttons report success in text.
    \item Color is never the only status indicator.
    \item Controls remain usable at 200\% zoom.
    \item Mobile input controls do not overflow.
    \item Motion is unnecessary for converter operation.
\end{itemize}

\subsection{UI planning prompt}

\begin{promptbox}{Plan the Interactive CU-Time Interface}
Using the approved engine API and the existing CU-Time page design,
create the implementation specification for the first interactive
converter release.

Do not edit files.

Specify:

1. Semantic form structure
2. Forward and reverse modes
3. Exact field labels and help text
4. Validation presentation
5. Result region and fields
6. Copy and clear behavior
7. aria-live behavior
8. Keyboard behavior
9. Responsive layout
10. Astro script loading and module imports
11. Error-state styling using existing tokens
12. Research classification and methodology disclosure
13. Deferred features
14. Manual test procedure
15. Acceptance criteria

Do not add React, Tailwind, or CDN calculation dependencies.
\end{promptbox}

\subsection{UI implementation prompt}

\begin{promptbox}{Integrate the Approved Converter with Astro}
Implement only the approved first-release interactive CU-Time
interface.

Preconditions:

- active branch is feature/cu-time-converter
- engine tests pass
- golden parity tests pass
- npm run build succeeds before UI changes

Requirements:

- update CUTimeConverter.astro
- connect the approved pure TypeScript engine
- use lightweight browser TypeScript
- preserve the surrounding cu-time.astro documentation
- preserve research classifications
- use existing design tokens
- provide visible labels and accessible errors
- announce results with aria-live
- preserve full precision as text
- include copy and clear controls
- support desktop and mobile layouts
- do not add React
- do not add Tailwind
- do not use CDN calculation libraries
- do not modify the reference HTML converter
- run npm test
- run npm run build
- report all changed files and manual test results

Do not commit, tag, push, merge, or deploy.
\end{promptbox}

\subsection{Local manual test}

\begin{commandbox}{Run the integrated page locally}
cd "$(git rev-parse --show-toplevel)/website" || exit 1

npm run dev
\end{commandbox}

Open:

\begin{codebox}{Local CU-Time route}
http://localhost:4321/CosmicUniversalismStatement/cu-time/
\end{codebox}

Test keyboard navigation, mobile width, all golden inputs, invalid inputs, copy,
clear, and browser-console output.

% ============================================================
\section{Part 11 - Acceptance, Review, Milestones, and Merge}
% ============================================================

\subsection{Parity report}

Before merge, create a report that maps every approved reference vector to the
new engine result.

\begin{longtable}{L{0.55in}L{1.45in}L{1.45in}L{0.9in}L{1.1in}}
\toprule
\textbf{ID} & \textbf{Reference} & \textbf{TypeScript} & \textbf{Result} & \textbf{Decision} \\
\midrule
 & & & Pass/Fail & \\
 & & & Pass/Fail & \\
 & & & Pass/Fail & \\
 & & & Pass/Fail & \\
\bottomrule
\end{longtable}

\subsection{Final acceptance checklist}

\begin{itemize}
    \item[\checkbox] Active work remained on \branchname{feature/cu-time-converter}.
    \item[\checkbox] \branchname{main} was not modified.
    \item[\checkbox] Original HTML reference converter is unchanged.
    \item[\checkbox] Reference hash matches the recorded hash.
    \item[\checkbox] Constants registry is approved.
    \item[\checkbox] Calendar specification is approved.
    \item[\checkbox] Forward specification is approved.
    \item[\checkbox] Reverse specification is approved.
    \item[\checkbox] Golden vectors are approved.
    \item[\checkbox] Decimal precision is preserved from strings.
    \item[\checkbox] Automated tests pass.
    \item[\checkbox] Golden parity tests pass.
    \item[\checkbox] Round-trip tests pass within documented tolerance.
    \item[\checkbox] Invalid dates are handled as approved.
    \item[\checkbox] BCE behavior is documented and tested.
    \item[\checkbox] Rollover behavior is documented and tested.
    \item[\checkbox] Converter is interactive at \texttt{\TargetRoute}.
    \item[\checkbox] Keyboard and mobile testing pass.
    \item[\checkbox] No console errors remain.
    \item[\checkbox] \texttt{npm run build} succeeds.
    \item[\checkbox] Working tree contains only understood files.
    \item[\checkbox] Documentation and current-state record are updated.
    \item[\checkbox] Explicit approval to commit has been given.
    \item[\checkbox] Explicit approval to merge has been given.
\end{itemize}

\subsection{Required final verification commands}

\begin{commandbox}{Final feature verification}
cd "$(git rev-parse --show-toplevel)" || exit 1

git branch --show-current
git status --short --branch

cd website || exit 1
npm test
npm run build
cd ..

git diff --check
git status --short --branch
\end{commandbox}

\subsection{Recommended milestone sequence}

\begin{tabularx}{\textwidth}{L{2.55in}Y}
\toprule
\textbf{Milestone} & \textbf{Meaning} \\
\midrule
\nolinkurl{website-v1.1-cu-time-foundation} & Existing UI and repository foundation \\
\nolinkurl{website-v1.1-cu-time-specification} & Constants, calendar, formulas, and vectors approved \\
\nolinkurl{website-v1.1-cu-time-engine} & TypeScript engine and automated tests approved \\
\nolinkurl{website-v1.1-cu-time-converter} & Interactive page, parity, accessibility, and build approved \\
\bottomrule
\end{tabularx}

Tags should be created only after explicit approval and should point to clean,
reviewed commits. The final converter release tag is preferably created after
the approved feature is merged into \branchname{website}.

\subsection{Commit planning}

Suggested logical commits, subject to the actual approved work:

\begin{codebox}{Suggested commit sequence}
Document CU-Time mathematical specification
Add precision-safe CU-Time constants registry
Add CU-Time calendar conversion module
Add CU-Time conversion engine
Add CU-Time golden vector tests
Integrate interactive CU-Time converter
Document CU-Time parity and acceptance
\end{codebox}

\subsection{Merge procedure}

\begin{enumerate}
    \item Confirm all acceptance checks are complete.
    \item Confirm the feature branch is pushed and clean.
    \item Review the complete branch diff against \branchname{website}.
    \item Obtain explicit approval to merge.
    \item Merge \branchname{feature/cu-time-converter} into \branchname{website}.
    \item Run the tests and production build on \branchname{website}.
    \item Push \branchname{website} only after verification.
    \item Confirm the website deployment workflow completes.
    \item Verify the public CU-Time route.
    \item Create the approved release tag.
    \item Update both implementation manuals and the current-state record.
\end{enumerate}

\begin{warningbox}
The deployment workflow is configured to run on pushes to
\branchname{website}, not the feature branch. Merging and pushing to
\branchname{website} may publish the converter. Treat merge approval as release
approval unless the workflow is intentionally disabled first.
\end{warningbox}

\subsection{Final review prompt}

\begin{promptbox}{Final CU-Time Converter Review}
Perform a final no-edit review of the CU-Time converter feature.

Compare feature/cu-time-converter against website.

Verify:

1. Branch and working-tree state
2. Original reference converter hash
3. Constants and precision representation
4. Calendar specification compliance
5. Forward and reverse specification compliance
6. Golden vector parity
7. Automated test results
8. Round-trip results
9. Validation behavior
10. Accessibility and responsive behavior
11. Target route and Astro base-path behavior
12. Dependency changes
13. Files outside the authorized scope
14. Build success
15. Documentation updates
16. Deployment implications
17. Remaining open questions

Report findings by severity.
Do not edit, commit, tag, push, merge, or deploy.
\end{promptbox}

\subsection{Final controlled statement}

\begin{tcolorbox}[
    enhanced,
    breakable,
    colback=CUNavy,
    colframe=CUGold,
    coltitle=CUNavy,
    colbacktitle=CUGoldLight,
    boxrule=1pt,
    arc=2mm,
    title=\textbf{Controlled Migration Rule}
]
\color{white}

The Cosmic Breath Time Converter must be migrated by preserving observable
reference behavior, documenting constants and calendar assumptions, capturing
golden vectors, implementing pure precision-safe TypeScript functions, proving
parity through automated tests, and only then connecting the engine to the
Astro interface.

The original HTML file remains unchanged during migration. Potential
mathematical or implementation corrections are documented and approved
independently; they are never introduced silently under the label of parity.

\end{tcolorbox}

% ============================================================
\appendix
\section{Appendix A - Current-State Record Template}
% ============================================================

After implementation begins, maintain a repository record such as:

\begin{codebox}{website/docs/cu-time/CU\_TIME\_CURRENT\_STATE.md}
# CU-Time Current State

Repository: willmaddock/CosmicUniversalismStatement
Protected branch: main
Stable website branch: website
Active feature branch: feature/cu-time-converter
Current HEAD: [UPDATE]
Current milestone: [UPDATE]
Current tags: [UPDATE]
Target route: /CosmicUniversalismStatement/cu-time/
Reference converter hash: [UPDATE]

## Completed
- [UPDATE]

## In Progress
- [UPDATE]

## Pending
- [UPDATE]

## Tests
- npm test: [RESULT]
- npm run build: [RESULT]

## Open Questions
- [UPDATE]

## Merge Status
- [UPDATE]
\end{codebox}

% ============================================================
\section{Appendix B - Phase Authorization Prompt}
% ============================================================

Use this prompt before every implementation phase.

\begin{promptbox}{Authorize One Controlled Phase}
We are authorizing one controlled CU-Time implementation phase.

Authorized phase:
[INSERT PHASE]

Approved specification:
[INSERT OR REFERENCE APPROVED SECTION]

Before editing:

1. Confirm branch feature/cu-time-converter.
2. Confirm working-tree status.
3. List files expected to change.
4. List dependencies expected to change.
5. State the tests that must pass.
6. State the completion gate.

During implementation:

- change only the authorized scope
- preserve the reference HTML file
- preserve exact Decimal literals
- do not silently correct behavior
- document all deviations

After implementation:

- run the authorized tests
- run npm run build
- report every changed file
- report warnings, failures, and open questions

Do not commit, tag, push, merge, or deploy.
\end{promptbox}

% ============================================================
\section{Appendix C - Manual Revision Record}
% ============================================================

\begin{tabularx}{\textwidth}{L{0.65in}L{1.2in}Y}
\toprule
\textbf{Version} & \textbf{Date} & \textbf{Change} \\
\midrule
1.0 & July 19, 2026 & Initial controlled converter migration manual based on the verified CU-Time foundation branch, repository inspection, and successful Astro production build. \\
\bottomrule
\end{tabularx}

\vfill

\begin{center}
\color{CUMuted}
\small
End of \ProjectTitle\ \DocumentTitle
\end{center}

\end{document}
